\documentclass[a4paper,11pt]{ctexart}
\usepackage{amsmath, amssymb, amsfonts}
\usepackage{geometry}
\geometry{left=1.5cm, right=1.5cm, top=2.0cm, bottom=2.0cm, headsep=0.3cm, footskip=1cm}
\usepackage{listings}
\usepackage{xcolor}
\usepackage{tikz}
\usepackage{pgfplots}
\usepackage{hyperref}
\usepackage{fancyhdr}
\usepackage{tcolorbox}
\tcbuselibrary{breakable, skins}
\usepackage{array}
\usepackage{booktabs}
\usepackage[shortlabels]{enumitem}
\usepackage{needspace}
\usepackage{multirow}

\AtBeginDocument{
  \hypersetup{colorlinks=true, linkcolor=blue, filecolor=magenta, urlcolor=cyan}
}

\setlist{nosep, itemsep=0.8ex, parsep=0.1em}
\setlist[itemize]{leftmargin=1.8em}
\setlist[enumerate]{leftmargin=1.8em}

% ===== 颜色定义 =====
\definecolor{codegreen}{rgb}{0,0.6,0}
\definecolor{codegray}{rgb}{0.5,0.5,0.5}
\definecolor{codepurple}{rgb}{0.58,0,0.82}
\definecolor{codebg}{rgb}{0.98,0.98,0.98}
\definecolor{tipsblue}{rgb}{0.85,0.93,1.0}
\definecolor{highlightyellow}{rgb}{1.0,0.97,0.8}

% ===== 代码块样式 =====
\lstdefinestyle{mystyle}{
    backgroundcolor=\color{codebg},
    commentstyle=\color{codegreen},
    keywordstyle=\color{magenta}\bfseries,
    numberstyle=\tiny\color{codegray},
    stringstyle=\color{codepurple},
    basicstyle=\ttfamily\small,
    breakatwhitespace=false,
    breaklines=true,
    captionpos=b,
    keepspaces=true,
    numbers=left,
    numbersep=6pt,
    showspaces=false,
    showstringspaces=false,
    showtabs=false,
    tabsize=2,
    language=Python,
    frame=single,
    framerule=0.5pt,
    rulecolor=\color{gray!40}
}
\lstset{style=mystyle}

% ===== 彩色框 =====
\newtcolorbox{problembox}[1][]{
    colback=white,
    colframe=blue!60!black,
    title=\textbf{#1},
    fonttitle=\bfseries\small,
    boxrule=1pt,
    arc=3pt, outer arc=3pt,
    coltitle=black,
    before skip=10pt, after skip=8pt
}

\newtcolorbox{solutionbox}[1][]{
    enhanced, breakable,
    colback=green!3!white,
    colframe=green!50!black,
    title=\textbf{#1},
    fonttitle=\bfseries\small,
    boxrule=1pt,
    arc=3pt, outer arc=3pt,
    coltitle=black,
    before skip=8pt, after skip=10pt
}

\newtcolorbox{metaphorbox}[1][]{
    colback=orange!5!white,
    colframe=orange!70!black,
    title=\textbf{#1},
    fonttitle=\bfseries\small,
    boxrule=1pt,
    arc=3pt, outer arc=3pt,
    coltitle=black,
    before skip=8pt, after skip=8pt
}

\newtcolorbox{beginnerbox}[1][]{
    colback=tipsblue,
    colframe=blue!50!black,
    title=\textbf{#1},
    fonttitle=\bfseries\small,
    boxrule=1pt,
    arc=3pt, outer arc=3pt,
    coltitle=black,
    before skip=8pt, after skip=8pt
}

\newtcolorbox{appbox}[1][]{
    colback=green!3!white,
    colframe=green!60!black,
    title=\textbf{#1},
    fonttitle=\bfseries\small,
    boxrule=1pt,
    arc=3pt, outer arc=3pt,
    coltitle=black,
    before skip=8pt, after skip=10pt
}

\newtcolorbox{keybox}[1][]{
    colback=highlightyellow,
    colframe=orange!80!black,
    title=\textbf{#1},
    fonttitle=\bfseries\small,
    boxrule=1pt,
    arc=3pt, outer arc=3pt,
    coltitle=black,
    before skip=8pt, after skip=8pt
}

% ===== 页眉页脚 =====
\pagestyle{fancy}
\fancyhf{}
\fancyhead[R]{深度学习-第9章}
\fancyhead[L]{RNN 循环神经网络 · 练习题}
\fancyfoot[C]{\thepage}
\addtolength{\headheight}{2pt}

\parskip=2pt plus 1pt minus 0.5pt
\linespread{1.35}
\raggedbottom
\widowpenalty=10000
\clubpenalty=10000

\title{\textbf{深度学习\\第9章\quad 循环神经网络（RNN）}\\[0.5em]
{\large\color{gray}——当神经网络需要"记住"过去，该怎么做？}}
\date{\today}

\begin{document}

\maketitle

% ===== 章节导读 =====
\begin{beginnerbox}[本章题目导览：你将挑战哪些核心问题？]
    CNN 擅长处理静态图像，但现实中还有大量\textbf{顺序有关}的数据：你说的每一句话、写的每一行代码、股票的每一天价格……这些都是"时序数据"，顺序一变，意思就变。本章 6 道题从零打通 RNN 的核心脉络：

    \begin{itemize}
        \item \textbf{Q9-01（辨析填空）}：NLP 三种方法横向对比——同义词词典 vs 基于计数 vs 基于推理，各自解决什么问题、有何缺陷？
        \item \textbf{Q9-02（概念理解）}：词嵌入 vs 独热向量——为什么独热向量表示词语很"蠢"？词嵌入怎么就聪明了？
        \item \textbf{Q9-03（计算题）}：RNN 隐藏状态手动推导——给定权重和当前输入，一步步算出 $h_t$，把公式变成具体数字。
        \item \textbf{Q9-04（对比分析）}：RNN vs 前馈网络——权重共享、时间展开，RNN 到底和普通神经网络有什么本质区别？
        \item \textbf{Q9-05（代码阅读）}：古诗生成模型代码追踪——\texttt{PoetryRNN} 的数据流形状逐层追踪，损失函数为何要 \texttt{transpose}？
        \item \textbf{Q9-06（开放分析）}：文本生成中的"温度"采样——为什么用 \texttt{multinomial} 而不用 \texttt{argmax}？温度参数如何控制创作风格？
    \end{itemize}
\end{beginnerbox}

\section{第 9 章\quad 循环神经网络（RNN）与自然语言处理}

在学 CNN 的时候，我们处理的都是"静态的"东西——一张图片，像素全部同时摆在那里，没有顺序关系。

但今天要处理的东西不一样了：

\begin{itemize}
    \item "我喜欢你"和"你喜欢我"——词完全一样，顺序不同，意思天差地别；
    \item 明天的气温，取决于今天、昨天、前天……
    \item 下一个字是什么，取决于你前面已经打了什么。
\end{itemize}

这类数据有一个共同特点：\textbf{顺序很重要，上下文很重要}。处理它们，就需要 RNN。

\begin{metaphorbox}[先建立一个总感觉：RNN 是什么]
    \textbf{把 RNN 想象成一个边看书边记笔记的人。}

    \begin{itemize}
        \item 普通神经网络（CNN、全连接）= 每次只看当前这一张卡片，上一张忘得干干净净；
        \item RNN = 每读到一个新词，会翻一眼自己之前的笔记，然后把新信息和旧笔记合并，更新成一份新笔记，传给下一步继续用。
    \end{itemize}

    RNN 最核心的能力就是四个字：\textbf{带着记忆往前走。}
\end{metaphorbox}


%% ============================================================
\Needspace{5\baselineskip}
\begin{problembox}[Q9-01\quad NLP 三种核心方法横向对比]
    \textbf{题目描述：}\\
    文档第 9.1 节介绍了计算机理解自然语言的三种方法：基于同义词词典、基于计数、基于推理。在做题之前，先来认识一下这三种思路的本质区别。

    \vspace{0.3cm}
    \textbf{问题：}
    \begin{enumerate}[(1)]
        \item 填写下表，将三种方法的关键特征补全：

        \begin{center}
        \begin{tabular}{p{3.2cm} p{3.5cm} p{4.5cm} p{2.8cm}}
            \toprule
            方法 & 核心思想 & 主要缺点 & 是否自动化 \\
            \midrule
            基于同义词词典 & 人工构建词语网络，计算词语\underline{\hspace{1.5cm}} & \underline{\hspace{3.5cm}}（写出2条） & \underline{\quad} \\[1.5em]
            基于计数 & 统计词频，用\underline{\hspace{1.5cm}}（矩阵名称）表示词向量 & 对所有词\underline{\hspace{2.5cm}}极高 & \underline{\quad} \\[1.5em]
            基于推理 & 根据\underline{\hspace{1.5cm}}预测当前位置的词 & 需要大量数据和计算资源 & \underline{\quad} \\
            \bottomrule
        \end{tabular}
        \end{center}

        \item "语料库"（corpus）是什么？基于计数的方法为什么需要它？
        \item 基于推理的方法和基于计数的方法，在"如何获得词向量"这件事上，根本区别是什么？
    \end{enumerate}

    \vspace{0.2cm}
    {\color{gray}\small\textit{来源溯源：[文档第 9 章 9.1 自然语言处理概述]}}
\end{problembox}

\begin{beginnerbox}[小白必读：为什么计算机理解语言这么难？]
    \begin{itemize}
        \item \textbf{编程语言 vs 自然语言}：编程语言语法严格，一个符号只有一个含义，计算机按规则解析即可。但汉语英语是"软语言"——同一个词在不同语境下含义可能完全不同（"苹果"可以是水果，也可以是手机品牌），新词不断涌现，规则无法穷举。
        \item \textbf{计算机只认数字}：不管哪种方法，最终目的都是把"词"变成"向量（数字列表）"，让计算机能够进行数学运算——计算两个词有多相似、预测下一个词是什么。
        \item \textbf{三种方法的进化逻辑}：词典方法靠人工→计数方法靠统计→推理方法靠神经网络。越往后，越自动化，越能处理大规模数据，但也越需要算力。
    \end{itemize}
\end{beginnerbox}

\begin{metaphorbox}[生活比喻：三种"认识一个词"的方式]
    假设你是一个刚学中文的外国人，想搞懂"秋"是什么意思：

    \begin{itemize}
        \item \textbf{词典法}：翻字典，发现"秋"=秋天、金秋、丰收……手工查好每个词的关系。但新词没收录怎么办？字典要一直更新，太累了。
        \item \textbf{计数法}：读了一万篇文章，发现"秋"总是和"风""叶""凉""稻"一起出现——靠这些邻居来猜它的意思。但要数遍所有文章里所有词两两共现，计算量爆炸。
        \item \textbf{推理法}：让神经网络做填空题："金风玉露一相\underline{\quad}"——靠这种上下文预测任务，网络自然就学会了每个词的"语义"，也就是词向量。这是现代 NLP 的核心思路。
    \end{itemize}
\end{metaphorbox}

\begin{solutionbox}[参考答案与详细解析]
    \textbf{(1) 填表：}

    \begin{center}
    \begin{tabular}{p{3.0cm} p{3.5cm} p{4.5cm} p{2.5cm}}
        \toprule
        方法 & 核心思想 & 主要缺点 & 是否自动化 \\
        \midrule
        基于同义词词典 & 人工构建词语网络，计算词语\textbf{相似度} & (1)需要人工逐个定义，费时费力；(2)新词维护成本极高；(3)表现力有限 & \textbf{否} \\[0.5em]
        基于计数 & 统计词频，用\textbf{共现矩阵}表示词向量 & 对所有词向量化的\textbf{计算复杂度}极高 & \textbf{是} \\[0.5em]
        基于推理 & 根据\textbf{上下文}预测当前位置的词 & 需要大量数据和计算资源 & \textbf{是} \\
        \bottomrule
    \end{tabular}
    \end{center}

    \textbf{(2) 语料库（corpus）的定义与作用：}

    语料库是\textbf{大量的文本数据集合}。基于计数的方法需要统计词频——即每个词和哪些词经常出现在一起（共现关系），这个统计必须建立在海量真实文本的基础上。语料库越大，统计出的词向量就越能反映真实的语言规律。没有语料库，就没有可统计的数据。

    \textbf{(3) 计数 vs 推理：获得词向量的根本区别：}

    \begin{keybox}[核心对比]
        \begin{itemize}
            \item \textbf{基于计数}：先统计整个语料库中的词频共现，\textbf{一次性算出}所有词的向量（共现矩阵的某列/某行）——这是统计方法，\textbf{不需要训练模型}。
            \item \textbf{基于推理}：设计一个神经网络，让它反复做"根据上下文猜词"的任务，\textbf{在训练过程中}逐渐学会词的向量表示——这是学习方法，\textbf{需要迭代优化模型}。
        \end{itemize}
        推理方法的词向量（也就是词嵌入）是训练神经网络时"顺带学出来"的副产品，同时也是 NLP 模型的第一层输入。Word2Vec、GloVe 这些著名的词向量技术，本质上都是基于推理思想的变体。
    \end{keybox}
\end{solutionbox}

\begin{appbox}[现实应用场景]
    \textbf{词向量让搜索引擎"懂"你的意思：}\\
    早期搜索引擎只做关键词匹配——你搜"苹果手机"，它只返回包含这两个字的页面。基于词嵌入的现代搜索引擎能理解语义：搜"苹果手机"，它知道"iPhone"和这四个字语义相近，也会返回相关结果。这背后，就是训练好的词嵌入模型把"苹果手机"和"iPhone"映射到了向量空间中相近的位置。Google 的 BERT、百度的 ERNIE 等大语言模型，本质上都在做更深层的"理解语义"工作，其基础正是词嵌入。
\end{appbox}


%% ============================================================
\Needspace{5\baselineskip}
\begin{problembox}[Q9-02\quad 词嵌入 vs 独热向量：谁更聪明？]
    \textbf{题目描述：}\\
    文档 9.2.1 节提到，独热向量（One-Hot Vector）虽然容易构建，但存在两个致命缺陷，因此现代 NLP 普遍使用词嵌入（Word Embedding）来表示词语。

    \vspace{0.2cm}
    \textbf{背景知识：}\\
    假设词表共有 $V=5$ 个词：\{我，喜欢，苹果，香蕉，手机\}。

    独热向量表示如下（每个词只有一个位置为 1，其余全为 0）：
    \[
    \text{我} = [1,0,0,0,0], \quad \text{喜欢} = [0,1,0,0,0], \quad \text{苹果} = [0,0,1,0,0]
    \]
    \[
    \text{香蕉} = [0,0,0,1,0], \quad \text{手机} = [0,0,0,0,1]
    \]

    词嵌入则是：把每个词映射到一个低维稠密向量（如 3 维），这些向量由训练得到：
    \[
    \text{苹果} \approx [0.8, 0.2, 0.1], \quad \text{香蕉} \approx [0.7, 0.3, 0.0], \quad \text{手机} \approx [0.1, 0.9, 0.8]
    \]

    \vspace{0.2cm}
    \textbf{问题：}
    \begin{enumerate}[(1)]
        \item 计算"苹果"和"香蕉"的独热向量余弦相似度；再计算"苹果"和"手机"的独热向量余弦相似度。结果说明了什么问题？
        \item 若词表大小为 $V=50000$（这是一个正常的中等词表大小），独热向量每个需要多少维？相比之下，如果用 $d=256$ 维的词嵌入，存储量之比大约是多少？
        \item 文档中说词嵌入矩阵（\texttt{nn.Embedding}）的参数形状是 $(\texttt{num\_embeddings}, \texttt{embedding\_dim})$，它在神经网络中是如何工作的？用"查表"比喻解释。
    \end{enumerate}

    \vspace{0.2cm}
    {\color{gray}\small\textit{来源溯源：[文档第 9 章 9.2.1 什么是词嵌入]}}
\end{problembox}

\begin{beginnerbox}[小白必读：余弦相似度是什么？]
    余弦相似度是衡量两个向量"方向有多接近"的指标，公式为：
    \[
    \cos(\vec{a}, \vec{b}) = \frac{\vec{a} \cdot \vec{b}}{|\vec{a}||\vec{b}|}
    \]
    \begin{itemize}
        \item 结果范围：$[-1, 1]$，越接近 1 说明越相似，等于 0 说明完全无关，等于 $-1$ 说明完全相反。
        \item \textbf{对独热向量}：任意两个不同词的独热向量点积为 0（因为没有同为 1 的位置），所以余弦相似度恒为 0——"苹果"和"香蕉"跟"苹果"和"手机"，在独热空间里\textbf{一样不相关}，完全无法区分"语义近"和"语义远"。
    \end{itemize}
\end{beginnerbox}

\begin{metaphorbox}[生活比喻：独热向量是穿着囚服的词，词嵌入是有身份证的词]
    \begin{itemize}
        \item \textbf{独热向量（囚服）}：每个词穿着一件只有自己编号的囚服，外表完全相同（都是一串 0 里有一个 1），除了编号外传达不了任何信息。问"苹果和香蕉像不像"？回答：一个是 3 号，一个是 4 号，一样的囚服，没区别。
        \item \textbf{词嵌入（身份证）}：每个词有一张丰富的身份证，上面写着它的各种"属性数值"（颜色倾向、可食用程度、科技感……）。苹果和香蕉的身份证数字很接近（都是水果），而苹果和手机的数字差很远（一个是食物，一个是电子产品）。
        \item \textbf{学习过程}：这张"身份证"不是人工填写的，而是神经网络在大量文本中自动学到的——它发现苹果和香蕉总在相似语境里出现，就自动给它们分配了相近的向量。
    \end{itemize}
\end{metaphorbox}

\begin{solutionbox}[参考答案与详细解析]
    \textbf{(1) 独热向量的余弦相似度：}

    "苹果" $=[0,0,1,0,0]$，"香蕉" $=[0,0,0,1,0]$，"手机" $=[0,0,0,0,1]$。

    任意两个不同词的独热向量点积：
    \[
    \text{苹果} \cdot \text{香蕉} = 0\times0 + 0\times0 + 1\times0 + 0\times1 + 0\times0 = 0
    \]
    \[
    \text{苹果} \cdot \text{手机} = 0
    \]

    所有独热向量模长均为 1，因此：
    \[
    \cos(\text{苹果}, \text{香蕉}) = \frac{0}{1 \times 1} = 0 \quad \cos(\text{苹果}, \text{手机}) = 0
    \]

    \textbf{结论}：苹果和香蕉（语义相近，都是水果）与苹果和手机（语义相远）的相似度\textbf{完全相同，都是 0}。独热向量无法区分语义近远，完全丢失了词语之间的语义关系。

    \vspace{0.5em}
    \textbf{(2) 存储量对比：}

    独热向量维度 $= V = 50000$，每个词占 50000 个数。\\
    词嵌入维度 $= d = 256$，每个词只占 256 个数。

    存储量之比：
    \[
    \frac{50000}{256} \approx 195 \text{ 倍}
    \]

    词嵌入的总参数量 $= V \times d = 50000 \times 256 = 1280$ 万个，虽然看起来多，但\textbf{每个词只用 256 个数表示}，比独热向量的 50000 维大幅压缩，且包含了语义信息。

    \vspace{0.5em}
    \textbf{(3) \texttt{nn.Embedding} 的工作原理——查表：}

    词嵌入矩阵形状为 $(V, d)$，可以理解成一本\textbf{"词向量字典"}：第 $i$ 行就是第 $i$ 个词的 $d$ 维向量。

    \begin{keybox}[工作流程：三步查表]
        \begin{enumerate}
            \item \textbf{分词 + 建词表}：把文本分成词，给每个词分配一个整数 ID（"我"=0，"苹果"=2……）
            \item \textbf{查表}：把词的 ID 当作行索引，直接取出嵌入矩阵对应的那一行向量。\texttt{embed(torch.tensor(2))} 就是取出第 2 行，也就是"苹果"的词向量。
            \item \textbf{梯度更新}：这个矩阵的每一行都是可学习参数。训练过程中，根据损失反向传播，不断调整每一行的数值，让语义相近的词向量越来越靠近。
        \end{enumerate}
    \end{keybox}

    \begin{lstlisting}[language=Python]
import torch, torch.nn as nn

# 词表大小5，词向量维度3
embed = nn.Embedding(num_embeddings=5, embedding_dim=3)

# "苹果"的 ID 是 2，直接取出第2行
apple_vec = embed(torch.tensor(2))
print(apple_vec)  # 输出 3 维向量，初始是随机值，训练后有语义

# 一次查多个词
batch_ids = torch.tensor([0, 2, 4])  # 我、苹果、手机
batch_vecs = embed(batch_ids)
print(batch_vecs.shape)  # torch.Size([3, 3])
    \end{lstlisting}
\end{solutionbox}

\begin{appbox}[现实应用场景]
    \textbf{Word2Vec：第一个让词有"方向感"的词嵌入：}\\
    2013 年 Google 发布的 Word2Vec，通过在大规模语料上训练"预测上下文"任务，学出了惊人的词向量性质：$\vec{\text{国王}} - \vec{\text{男}} + \vec{\text{女}} \approx \vec{\text{女王}}$。这说明词向量空间里自发出现了"性别方向"和"王室方向"这样的语义轴。现代大语言模型（GPT、BERT）的 embedding 层依然继承了这一思想，只是用 Transformer 代替了简单的神经网络来训练更丰富的词语表示。
\end{appbox}


%% ============================================================
\Needspace{5\baselineskip}
\begin{problembox}[Q9-03\quad RNN 隐藏状态手动推导]
    \textbf{题目描述：}\\
    RNN 层的核心公式是：
    \begin{equation}
        h_t = \tanh(h_{t-1} W_h + x_t W_x + b)
    \end{equation}
    其中 $x_t$ 是当前输入，$h_{t-1}$ 是上一时刻的隐藏状态，$W_x$、$W_h$ 是权重矩阵，$b$ 是偏置。

    \textbf{注意：}文档中标记符号为 $W_{xh}$ 对应 $W_x$，$W_{hh}$ 对应 $W_h$。

    \vspace{0.2cm}
    \textbf{具体数值：}\\
    已知某 RNN 层中（$\tanh$ 激活\textbf{之前}的线性部分）：

    \[
    x_t = \begin{bmatrix} 2.0 \\ -1.0 \end{bmatrix}, \quad
    h_{t-1} = \begin{bmatrix} 0.0 \\ 0.0 \end{bmatrix}, \quad
    W_x = \begin{bmatrix} 0.5 & -0.5 \\ 0.5 & 0.5 \end{bmatrix}, \quad
    W_h = \begin{bmatrix} 1.0 & 0.0 \\ 0.0 & 1.0 \end{bmatrix}, \quad
    b = \begin{bmatrix} 0.0 \\ 0.0 \end{bmatrix}
    \]

    \vspace{0.2cm}
    \textbf{问题：}
    \begin{enumerate}[(1)]
        \item 计算线性组合 $a_t = h_{t-1}W_h + x_t W_x + b$，写出完整推导过程。
        \item 在 $a_t$ 上应用 $\tanh$ 激活函数，得出最终隐藏状态 $h_t$。\\（提示：$\tanh(1.5) \approx 0.905$，$\tanh(-0.5) \approx -0.462$）
        \item 如果这是一首古诗的第一个字（$t=1$），隐藏状态 $h_0$ 应该初始化为多少？代码里是如何传入初始隐藏状态的？
    \end{enumerate}

    \vspace{0.2cm}
    {\color{gray}\small\textit{来源溯源：[文档第 9 章 9.3.1 RNN 层介绍]}}
\end{problembox}

\begin{beginnerbox}[小白必读：$h_t$ 的两个"来源"是什么？]
    每一步 RNN 的输入有两个：
    \begin{itemize}
        \item \textbf{当前词 $x_t$}：这一时刻刚刚输入的信息（比如古诗第 $t$ 个字的词向量）；
        \item \textbf{历史记忆 $h_{t-1}$}：上一时刻积累下来的"状态摘要"，包含了前面所有时刻的信息痕迹。
    \end{itemize}
    这两者分别经过各自的权重矩阵"加工"，再合在一起，过一次 $\tanh$ 压缩，就得到了新的记忆 $h_t$。

    \textbf{为什么用 $\tanh$ 而不用 ReLU？}\\
    $\tanh$ 的输出范围是 $(-1, 1)$。RNN 会把 $h_{t-1}$ 一次次地乘以 $W_h$ 往前传，如果不压缩，数值会越乘越大（或越乘越小），导致"梯度爆炸"或"梯度消失"。$\tanh$ 把每一步的输出都限制在 $(-1,1)$ 之间，起到"稳压器"的作用。
\end{beginnerbox}

\begin{metaphorbox}[生活比喻：每一步都在"消化新闻+更新认知"]
    把 RNN 的每一步想象成一个人在看新闻滚动条：

    \begin{itemize}
        \item $h_{t-1}$（旧认知）：他到上一条新闻为止，脑子里对世界的整体判断；
        \item $x_t$（当前新闻）：刚收到的这一条新消息；
        \item $W_h$（旧认知的权重）：他有多信任自己以往的判断；
        \item $W_x$（新消息的权重）：他有多重视当前这条新消息；
        \item $\tanh$（情绪稳定器）：不管收到多夸张的新闻，他的"认知更新量"都被控制在合理范围内，不会暴跳如雷（梯度爆炸）也不会完全无感（梯度消失）。
    \end{itemize}
\end{metaphorbox}

\begin{solutionbox}[参考答案与详细解析]
    \textbf{(1) 计算线性组合 $a_t$：}

    由于 $h_{t-1} = [0, 0]^\top$，所以历史记忆这一项贡献为零：
    \[
    h_{t-1} W_h = \begin{bmatrix} 0 & 0 \end{bmatrix}
    \begin{bmatrix} 1 & 0 \\ 0 & 1 \end{bmatrix}
    = \begin{bmatrix} 0 & 0 \end{bmatrix}
    \]

    当前输入的贡献（注意 $x_t$ 是列向量，与 $W_x$ 相乘为行向量，或按文档约定写为行乘矩阵）：

    这里我们采用行向量左乘矩阵的约定，$x_t^\top = [2.0, -1.0]$：
    \[
    x_t^\top W_x =
    \begin{bmatrix} 2.0 & -1.0 \end{bmatrix}
    \begin{bmatrix} 0.5 & -0.5 \\ 0.5 & 0.5 \end{bmatrix}
    =
    \begin{bmatrix}
        2.0 \times 0.5 + (-1.0) \times 0.5 \\
        2.0 \times (-0.5) + (-1.0) \times 0.5
    \end{bmatrix}^\top
    =
    \begin{bmatrix} 1.5 & -1.5 \end{bmatrix}
    \]

    等等，让我们重新按标准列向量约定来算（矩阵左乘列向量）：

    \[
    W_x \cdot x_t =
    \begin{bmatrix} 0.5 & -0.5 \\ 0.5 & 0.5 \end{bmatrix}
    \begin{bmatrix} 2.0 \\ -1.0 \end{bmatrix}
    =
    \begin{bmatrix}
        0.5 \times 2.0 + (-0.5) \times (-1.0) \\
        0.5 \times 2.0 + 0.5 \times (-1.0)
    \end{bmatrix}
    =
    \begin{bmatrix} 1.0 + 0.5 \\ 1.0 - 0.5 \end{bmatrix}
    =
    \begin{bmatrix} 1.5 \\ 0.5 \end{bmatrix}
    \]

    加上偏置 $b = [0, 0]^\top$，最终：
    \[
    \boxed{a_t = \begin{bmatrix} 1.5 \\ 0.5 \end{bmatrix}}
    \]

    \vspace{0.5em}
    \textbf{(2) 应用 $\tanh$ 激活，得到 $h_t$：}
    \[
    h_t = \tanh(a_t) =
    \begin{bmatrix} \tanh(1.5) \\ \tanh(0.5) \end{bmatrix}
    \approx
    \begin{bmatrix} 0.905 \\ 0.462 \end{bmatrix}
    \]

    \[
    \boxed{h_t \approx \begin{bmatrix} 0.905 \\ 0.462 \end{bmatrix}}
    \]

    \textbf{直觉解读：}第一步（$t=1$）只有当前词的贡献，没有历史记忆。$h_1$ 完全由第一个词决定，之后 $h_1$ 会作为 $h_{t-1}$ 传给第二步，把第一个词的信息带进去。

    \vspace{0.5em}
    \textbf{(3) 初始隐藏状态 $h_0$ 的初始化：}

    通常将 $h_0$ 初始化为\textbf{全零向量}，表示"序列开始前没有任何历史记忆"。

    在古诗生成代码中，对应的写法是 \texttt{hidden = None}，PyTorch 的 \texttt{nn.RNN} 接收到 \texttt{hx=None} 时会自动将初始状态设为全零：

    \begin{lstlisting}[language=Python]
# 古诗生成代码中的初始化片段（文档 9.4.5）
hidden = None  # 初始化隐状态为 None，等价于全零

# 第一步推理
input = torch.LongTensor([[start_token]]).to(device)
output, hidden = model(input, hidden)
# 此后 hidden 被更新，下一步可直接传入
output, hidden = model(next_input, hidden)
    \end{lstlisting}

    \begin{keybox}[形状提示：RNN API 的输入输出张量形状]
        \begin{itemize}
            \item \texttt{input}：形状 $[\text{seq\_len}, \text{batch\_size}, \text{input\_size}]$
            \item \texttt{hx}（初始隐状态）：形状 $[\text{num\_layers}, \text{batch\_size}, \text{hidden\_size}]$
            \item \texttt{output}：形状 $[\text{seq\_len}, \text{batch\_size}, \text{hidden\_size}]$
            \item \texttt{hn}（最终隐状态）：形状 $[\text{num\_layers}, \text{batch\_size}, \text{hidden\_size}]$
        \end{itemize}
    \end{keybox}
\end{solutionbox}


%% ============================================================
\Needspace{5\baselineskip}
\begin{problembox}[Q9-04\quad RNN 与前馈网络的本质区别]
    \textbf{题目描述：}\\
    文档 9.3.1 节指出，前馈网络的信号传播是\textbf{单向}的，而 RNN 的信号可以在层内\textbf{循环}。这两种架构在处理序列数据时有本质的区别。

    \vspace{0.2cm}
    \textbf{问题：}
    \begin{enumerate}[(1)]
        \item 判断题：下列说法正确吗？请逐条判断并给出理由。
        \begin{enumerate}[a.]
            \item "RNN 的 $W_h$ 在处理序列第 1 个词和第 20 个词时是两套不同的权重矩阵。"
            \item "前馈网络也能处理文本，只需要把句子所有词的独热向量拼起来作为输入即可。"
            \item "RNN 的输出 $h_t$ 只包含当前时刻 $x_t$ 的信息，不包含历史信息。"
        \end{enumerate}
        \item 文档说"将神经网络的层数加深，可以更有效地提取层次信息"。在 \texttt{nn.RNN} 中，参数 \texttt{num\_layers=2} 代表什么？第一个隐藏层的输出如何传给第二个隐藏层？
        \item 用下面的 PyTorch 代码，解释为什么 \texttt{output.shape} 是 \texttt{[1, 3, 16]} 而不是 \texttt{[1, 3, 8]}（8 是输入维度）？

        \begin{lstlisting}[language=Python]
import torch
rnn = torch.nn.RNN(input_size=8, hidden_size=16, num_layers=2)
input = torch.rand(1, 3, 8)
hx = torch.randn(2, 3, 16)
output, hn = rnn(input=input, hx=hx)
print(output.shape)   # torch.Size([1, 3, 16])
print(hn.shape)       # torch.Size([2, 3, 16])
        \end{lstlisting}
    \end{enumerate}

    \vspace{0.2cm}
    {\color{gray}\small\textit{来源溯源：[文档第 9 章 9.3 循环网络层]}}
\end{problembox}

\begin{beginnerbox}[小白必读：什么是"时间展开"？为什么说 RNN 有环路？]
    \begin{itemize}
        \item \textbf{环路的意思}：RNN 单元的输出 $h_t$ 会被"循环"送回到自身，作为下一个时刻的输入之一。这就是那个环路——$h_{t-1} \to \text{RNN} \to h_t \to \text{RNN} \to h_{t+1} \to \cdots$
        \item \textbf{时间展开}：处理一个长度为 $T$ 的序列时，你可以把这个 RNN 单元"展开"成 $T$ 个相连的副本，每个副本处理一个时刻，相邻副本之间通过 $h_t$ 传递信息。展开后看起来就像一个很深的前馈网络，但\textbf{所有副本共享同一套权重}。
        \item \textbf{权重共享 = 不管序列多长，参数总量固定}：不管句子有 10 个词还是 100 个词，RNN 的参数量不变，因为同一套权重重复使用了 10 次或 100 次。这和全连接网络完全不同——全连接网络如果输入从 10 维变成 100 维，参数量会乘以 10。
    \end{itemize}
\end{beginnerbox}

\begin{solutionbox}[参考答案与详细解析]
    \textbf{(1) 判断题逐条解析：}

    \begin{enumerate}[a.]
        \item \textbf{错误。}RNN 的核心特点正是\textbf{权重共享}——$W_h$、$W_x$、$b$ 在处理整个序列的每一个时间步时\textbf{完全相同}。序列的不同位置用的是同一套权重，只是传入的 $h_{t-1}$ 不同（携带了不同的历史信息），所以输出也不同。这正是 RNN 能以固定参数量处理任意长度序列的原因。

        \item \textbf{有一定道理，但有严重缺陷。}把所有词拼起来确实可以输入全连接层，但存在两大问题：(1)如果句子长度变化，输入维度就变化，全连接层无法处理不定长输入；(2)输入维度过大（一个词独热向量 $V$ 维，20 个词就是 $20V$ 维），参数量爆炸；(3)\textbf{顺序信息被丢弃}——拼接本身没有记录哪个词在前哪个词在后（当然可以靠位置编码补救，但这已经是 Transformer 的思路了）。

        \item \textbf{错误。}$h_t$ 不只包含当前时刻的信息，还包含\textbf{所有历史时刻的信息}（通过 $h_{t-1}$ 一层层积累传递下来）。这正是 RNN 的核心价值所在——$h_t$ 是一个"历史摘要+当前输入"的综合状态。
    \end{enumerate}

    \vspace{0.5em}
    \textbf{(2) \texttt{num\_layers=2} 的含义：}

    \texttt{num\_layers=2} 表示叠加了两层 RNN。第一层以原始输入序列 $x_t$ 为输入，输出隐藏序列 $h_t^{(1)}$；第二层以第一层的输出 $h_t^{(1)}$ 为输入，输出最终隐藏序列 $h_t^{(2)}$。\textbf{每一层都有独立的 $W_x$、$W_h$、$b$}。最终的 \texttt{output} 是第二层（最后一层）的输出，\texttt{hn} 是每一层最后时刻的隐藏状态。

    \vspace{0.5em}
    \textbf{(3) 为什么 \texttt{output.shape = [1, 3, 16]} 而不是 \texttt{[1, 3, 8]}：}

    \begin{keybox}[三个维度的含义]
        \begin{itemize}
            \item 第一维 \texttt{1}：序列长度 \texttt{seq\_len=1}（只处理了 1 个时刻的输入）
            \item 第二维 \texttt{3}：批量大小 \texttt{batch\_size=3}（同时处理 3 个样本）
            \item 第三维 \texttt{16}：\textbf{隐藏层大小 \texttt{hidden\_size=16}}，而不是输入大小 \texttt{input\_size=8}
        \end{itemize}
        输出维度是 \texttt{hidden\_size}，原因：RNN 的输出就是每个时刻的隐藏状态 $h_t$，而 $h_t$ 的维度是 \texttt{hidden\_size}，和输入维度 \texttt{input\_size} 无关。就像全连接层 \texttt{Linear(8, 16)} 的输出是 16 维，不是 8 维——输出维度由网络结构决定，不由输入维度决定。
    \end{keybox}
\end{solutionbox}

\begin{appbox}[现实应用场景]
    \textbf{多层 RNN 在语音识别中的应用：}\\
    深层双向 RNN（Bidirectional RNN）是 2014 年前语音识别的主流架构。第一层 RNN 学习声学特征的低级时序模式（音素级别），第二层学习更高级的语言结构（音节、词级别）。Google 的早期语音识别系统 DeepSpeech 就用了多层 RNN。不过从 2020 年起，Transformer 架构逐渐在语音识别中取代 RNN，主要原因是 Transformer 可以\textbf{并行训练}（RNN 必须逐步计算，无法并行），训练速度快得多。
\end{appbox}


%% ============================================================
\Needspace{5\baselineskip}
\begin{problembox}[Q9-05\quad 古诗生成代码追踪：数据流形状分析]
    \textbf{题目描述：}\\
    阅读文档 9.4 节的古诗生成案例代码，回答以下问题。

    \begin{lstlisting}[language=Python]
class PoetryRNN(nn.Module):
    def __init__(self, vocab_size, embedding_dim=128, hidden_size=256,
                 num_layers=1):
        super().__init__()
        self.embed = nn.Embedding(num_embeddings=vocab_size,
                                  embedding_dim=embedding_dim)
        self.rnn = nn.RNN(input_size=embedding_dim,
                          hidden_size=hidden_size, num_layers=num_layers)
        self.linear = nn.Linear(in_features=hidden_size,
                                out_features=vocab_size)

    def forward(self, input, hx=None):
        embed = self.embed(input)       # 第1步
        output, hidden = self.rnn(embed, hx)   # 第2步
        output = self.linear(output)    # 第3步
        return output, hidden

model = PoetryRNN(len(vocab), 256, 512, 2).to(device)
    \end{lstlisting}

    训练时，输入 \texttt{x} 的形状为 \texttt{[batch\_size=32, seq\_len=24]}（32 首诗，每首取 24 个字）。

    \vspace{0.2cm}
    \textbf{问题：}
    \begin{enumerate}[(1)]
        \item 假设 \texttt{vocab\_size=3000}，按照模型构造参数（\texttt{embedding\_dim=256, hidden\_size=512, num\_layers=2}），逐步追踪以下四个张量的形状：
        \begin{enumerate}[a.]
            \item \texttt{x}（模型输入）：\underline{\hspace{4cm}}
            \item \texttt{embed}（词嵌入层输出）：\underline{\hspace{4cm}}
            \item \texttt{output}（RNN 层输出）：\underline{\hspace{4cm}}
            \item \texttt{output}（线性层输出）：\underline{\hspace{4cm}}
        \end{enumerate}

        \item 训练代码中有这样一行：
        \begin{lstlisting}[language=Python]
loss_value = loss(output.transpose(1, 2), y)
        \end{lstlisting}
        为什么要 \texttt{transpose(1, 2)}？\texttt{CrossEntropyLoss} 期望的输入形状是什么？

        \item 训练时的目标序列 \texttt{y} 和输入序列 \texttt{x} 的关系是什么？为什么要这样设计？
    \end{enumerate}

    \vspace{0.2cm}
    {\color{gray}\small\textit{来源溯源：[文档第 9 章 9.4.2 自定义 Dataset / 9.4.3 搭建模型 / 9.4.4 模型训练]}}
\end{problembox}

\begin{beginnerbox}[小白必读：RNN 为什么要先过词嵌入再过 RNN？]
    \begin{itemize}
        \item \textbf{输入流水线}：文字 $\to$ 整数 ID $\to$ 词嵌入向量 $\to$ RNN $\to$ 预测下一个词的概率。
        \item \textbf{为什么不能直接把整数 ID 送进 RNN？}整数 ID（比如"月"=1234，"花"=567）是人工分配的编号，没有任何语义含义，1234 不比 567 更"月亮"。把这些编号直接送进 RNN，模型会把 ID 值当成数量关系来处理，这毫无意义。词嵌入的作用就是把无意义的整数变成有语义含义的浮点向量。
        \item \textbf{为什么要 \texttt{nn.RNN} 之后再接 \texttt{nn.Linear}？}RNN 的输出 $h_t$ 维度是 \texttt{hidden\_size=512}，但我们最终要预测的是"下一个字是词表中哪个字"，输出应该是 \texttt{vocab\_size=3000} 维的概率分布。\texttt{nn.Linear(512, 3000)} 就是把隐藏状态"解码"成词表上的 logits（未归一化概率）。
    \end{itemize}
\end{beginnerbox}

\begin{solutionbox}[参考答案与详细解析]
    \textbf{(1) 逐层形状追踪：}

    注意：\texttt{nn.RNN} 期望输入形状为 \texttt{[seq\_len, batch\_size, input\_size]}，但 \texttt{x} 的形状是 \texttt{[batch\_size, seq\_len]}，需要经过词嵌入和必要的维度调整。我们先追踪各层直接的输出形状：

    \begin{center}
    \begin{tabular}{clll}
        \toprule
        步骤 & 操作 & 输出形状 & 说明 \\
        \midrule
        输入 & \texttt{x} & \texttt{[32, 24]} & 32首诗，每首24字（整数ID） \\
        第1步 & \texttt{embed(x)} & \texttt{[32, 24, 256]} & 每个字ID变成256维词向量 \\
        （转置） & 调整维度喂给RNN & \texttt{[24, 32, 256]} & seq\_len在前，batch在中 \\
        第2步 & \texttt{rnn(embed)} & \texttt{[24, 32, 512]} & 每个时刻输出512维隐藏状态 \\
        第3步 & \texttt{linear(output)} & \texttt{[24, 32, 3000]} & 每个时刻预测3000个字的得分 \\
        \bottomrule
    \end{tabular}
    \end{center}

    \textbf{(2) 为什么要 \texttt{transpose(1, 2)}：}

    \texttt{CrossEntropyLoss} 对多类别分类的期望输入形状为：
    \[
    \texttt{[batch\_size, num\_classes, ...]}
    \]
    即类别维度必须是\textbf{第二个维度}（维度索引 1）。

    而 RNN 输出经过 \texttt{linear} 后形状是 \texttt{[seq\_len=24, batch\_size=32, vocab\_size=3000]}，如果直接换成批次在前，则是 \texttt{[32, 24, 3000]}——此时类别维度（3000）在第三个位置（索引 2），不满足 \texttt{CrossEntropyLoss} 的要求。

    \texttt{output.transpose(1, 2)} 交换第 1、2 维，把 \texttt{[32, 24, 3000]} 变成 \texttt{[32, 3000, 24]}，类别维度就到了索引 1，符合要求。

    \begin{keybox}[CrossEntropyLoss 期望的输入形状]
        \begin{itemize}
            \item \texttt{input}（预测）：形状 \texttt{[N, C]} 或 \texttt{[N, C, d1, d2, ...]}，$C$ 是类别数
            \item \texttt{target}（标签）：形状 \texttt{[N]} 或 \texttt{[N, d1, d2, ...]}，存储类别整数索引
            \item 本案例：\texttt{output.transpose(1,2)} 形状为 \texttt{[32, 3000, 24]}，\texttt{y} 形状为 \texttt{[32, 24]}，完全匹配
        \end{itemize}
    \end{keybox}

    \textbf{(3) 目标序列 \texttt{y} 和输入序列 \texttt{x} 的关系：}

    从 \texttt{PoetryDataset} 的代码可以看出：
    \begin{lstlisting}[language=Python]
self.data.append((seq[i : i + self.seq_len],        # x
                  seq[i + 1 : i + 1 + self.seq_len])) # y
    \end{lstlisting}

    \texttt{y} 比 \texttt{x}\textbf{整体向右偏移了 1 个字}——即 $y_t = x_{t+1}$。

    \textbf{这叫"下一个字预测"（Next Token Prediction）任务：}模型在看到 $x_t$ 时，要预测下一个字 $y_t = x_{t+1}$。训练时，已经知道了正确答案（$y$），通过交叉熵损失衡量预测有多准，再反向传播更新权重。这是语言模型的标准训练方式，GPT 系列大模型也是这样训练的，只是规模大了几千倍。
\end{solutionbox}


%% ============================================================
\Needspace{5\baselineskip}
\begin{problembox}[Q9-06\quad 文本生成中的采样策略：argmax vs multinomial]
    \textbf{题目描述：}\\
    文档 9.4.5 节的古诗生成代码中，生成下一个字的方式是：
    \begin{lstlisting}[language=Python]
prob = torch.softmax(output[0, 0], dim=-1)   # 计算概率分布
next_token = torch.multinomial(prob, 1)       # 从概率分布中随机采样
    \end{lstlisting}

    注意到这里用的是 \texttt{multinomial}（多项式采样），而不是 \texttt{argmax}（直接取概率最大的词）。

    \vspace{0.2cm}
    \textbf{问题：}
    \begin{enumerate}[(1)]
        \item 如果改用 \texttt{argmax} 策略，即每次都选概率最大的词，会产生什么问题？为什么古诗生成不适合这样做？
        \item \texttt{torch.softmax} 把 RNN 输出的 logits（原始得分）变成了概率分布，具体做了什么数学变换？如果词表里有 3 个词，logits 为 $[3.0, 1.0, 0.5]$，计算其 softmax 概率。（结果保留两位小数）
        \item 在文本生成中，有一种叫"温度（Temperature）"的参数，可以控制生成文本的"创意程度"：
        \[
        \text{prob} = \text{softmax}\!\left(\frac{\text{logits}}{T}\right)
        \]
        当 $T < 1$ 时，生成更保守（倾向于高概率词）；当 $T > 1$ 时，生成更随机（低概率词也有机会被选中）。请解释其背后的数学原因，并说明在古诗生成场景下，$T$ 选大一点还是小一点更好。
    \end{enumerate}

    \vspace{0.2cm}
    {\color{gray}\small\textit{来源溯源：[文档第 9 章 9.4.5 测试]}}
\end{problembox}

\begin{beginnerbox}[小白必读：Softmax 是什么？为什么必须先算 softmax 再采样？]
    \begin{itemize}
        \item \textbf{Softmax 的作用}：把任意实数向量（logits）变成\textbf{各元素之和为 1 的概率向量}。公式为 $p_i = e^{z_i} / \sum_j e^{z_j}$。无论输入多大多小，输出都是 $[0,1]$ 区间内的合法概率。
        \item \textbf{为什么要先 softmax 再采样？}\texttt{torch.multinomial} 需要概率分布作为输入，而 RNN 输出的 logits 可以是任意实数（可以是负数），不能直接当概率用。
        \item \textbf{argmax vs 采样的差别}：argmax 是确定性的，同一个模型在同样上下文下永远输出同样的词，一直重复。采样是随机的，带来多样性和"创意感"，非常适合诗歌这种需要变化的生成任务。
    \end{itemize}
\end{beginnerbox}

\begin{metaphorbox}[生活比喻：点菜时的"纠结程度"就是温度]
    你去餐厅点菜，菜单上有三道菜，你分别打了 $[3.0, 1.0, 0.5]$ 分：

    \begin{itemize}
        \item \textbf{argmax 策略（极度理性）}：每次都点评分最高的那道，从不尝试别的。时间长了，永远只吃那一道菜——模型也是如此，会产生极度单调重复的文本。
        \item \textbf{低温采样 $T=0.5$（保守型食客）}：评分差距被放大，高分菜的优势更明显，你大概率还是点高分菜，但偶尔也会鼓起勇气试试第二名。
        \item \textbf{高温采样 $T=2.0$（冒险型食客）}：评分差距被压缩，三道菜的概率趋于均匀，你可能随机点任何一道——模型生成的文本更有创意，但也可能产生"不知所云"的内容。
        \item \textbf{古诗生成的选择}：既要押韵、有章法（不能太随机），又要有一定新意（不能只会重复）。一般 $T \in [0.8, 1.2]$ 是较好的平衡点——稍微保守，但保留一定的随机性。
    \end{itemize}
\end{metaphorbox}

\begin{solutionbox}[参考答案与详细解析]
    \textbf{(1) argmax 策略的问题：}

    \textbf{贪心搜索的致命缺陷——重复生成}。如果每步都取概率最高的词，模型会陷入"循环陷阱"：比如生成到"春风吹又生"之后，下一步概率最高的词可能又是"春"，然后再是"风"……无限循环，无法生成有新意的内容。

    更根本的问题是：\textbf{局部最优不等于全局最优}。每步选最可能的词，不代表整首诗最可能。就像下棋，每步都走局部最好的一步，不一定赢得了整局。

    \vspace{0.5em}
    \textbf{(2) Softmax 计算：}

    Softmax 公式：$p_i = \dfrac{e^{z_i}}{\sum_{j} e^{z_j}}$

    给定 logits $= [3.0, 1.0, 0.5]$：

    \[
    e^{3.0} \approx 20.09, \quad e^{1.0} \approx 2.718, \quad e^{0.5} \approx 1.649
    \]
    \[
    \text{总和} = 20.09 + 2.718 + 1.649 = 24.457
    \]
    \[
    p_1 = \frac{20.09}{24.457} \approx 0.82, \quad
    p_2 = \frac{2.718}{24.457} \approx 0.11, \quad
    p_3 = \frac{1.649}{24.457} \approx 0.07
    \]

    \[
    \boxed{\text{softmax}([3.0, 1.0, 0.5]) \approx [0.82,\ 0.11,\ 0.07]}
    \]
    三者之和 $= 1.00$，满足概率分布要求。第一个词被选中的概率高达 82\%，体现了 softmax 对最大值的"放大"效果。

    \vspace{0.5em}
    \textbf{(3) 温度参数的数学原理：}

    \textbf{当 $T < 1$ 时（如 $T=0.5$）：}等效于把 logits 放大（$\text{logits}/0.5 = 2\times\text{logits}$），原来差距 2 的两个值变成差距 4。Softmax 对差距更敏感，高分值的概率被进一步放大，分布更"尖锐"（峰值更高），模型更倾向于选最高概率的词。

    \textbf{当 $T > 1$ 时（如 $T=2$）：}等效于把 logits 缩小（$\text{logits}/2$），原来差距 2 的两个值变成差距 1。Softmax 对差距不敏感，概率分布趋于均匀，各个词都有机会被选中，文本多样性增加。

    \textbf{当 $T \to 0$：}退化为 argmax（完全确定性）。

    \textbf{当 $T \to \infty$：}退化为均匀随机选词（完全随机）。

    \begin{keybox}[古诗生成中温度的实用建议]
        \begin{itemize}
            \item \textbf{$T < 0.7$}：输出较工整，但容易重复，缺乏新意
            \item \textbf{$T = 0.8 \sim 1.0$}：平衡模式，通常产生质量较好的古诗
            \item \textbf{$T > 1.2$}：创意较强但可能产生不押韵、无逻辑的输出
            \item \textbf{实践小技巧}：可以先用低温生成几个候选，再从中挑选最自然的一首
        \end{itemize}
    \end{keybox}

    \begin{lstlisting}[language=Python]
# 加入温度参数的生成代码（对原代码的改进版）
T = 0.9  # 温度参数，可以调节
prob = torch.softmax(output[0, 0] / T, dim=-1)  # 除以温度再 softmax
next_token = torch.multinomial(prob, 1)
    \end{lstlisting}
\end{solutionbox}

\begin{appbox}[现实应用场景]
    \textbf{ChatGPT 的 temperature 参数——你每天都在用的采样策略：}\\
    ChatGPT API 中有一个 \texttt{temperature} 参数（默认 1.0，范围 0--2），本质上就是本题讨论的温度采样。写代码时建议 $T=0.2$（更精确，避免随机错误）；写诗歌、故事时建议 $T=1.2$（更有创意）；需要稳定可复现结果时设 $T=0$（等价于 argmax）。OpenAI 还提供了 \texttt{top\_p}（核采样）参数作为温度的替代方案——只从累积概率超过 $p$ 的词中采样，也是控制多样性与准确性平衡的常用手段。
\end{appbox}

% ===== 本章小结 =====
\section{本章小结：把公式化成脑中的画面}

学完这 6 道题，如果还觉得"概念飘在空中"，试试把整章内容浓缩成这几句画面：

\begin{itemize}
    \item \textbf{NLP 三步进化}：人工词典（查字典）→ 统计共现（数邻居）→ 神经推理（做填空题）。
    \item \textbf{词嵌入}：把每个词变成一张"语义身份证"，训练完了语义相近的词向量也相近。
    \item \textbf{RNN 是什么}：边读序列边更新笔记，笔记（$h_t$）传给下一步继续用。
    \item \textbf{权重共享}：不管序列多长，始终用同一套规则处理每个时刻——这是 RNN 处理可变长序列的秘诀。
    \item \textbf{古诗生成的本质}：学会"已知前 $t$ 个字，预测第 $t+1$ 个字"，然后把预测结果不断喂回模型，一字一字地生成。
    \item \textbf{采样 vs argmax}：argmax 太死板，采样带来多样性，温度控制创意程度。
\end{itemize}

\begin{metaphorbox}[给零基础读者的一句总比喻]
    \textbf{如果把序列建模比作写一首歌：}
    \begin{itemize}
        \item 前馈网络：只看当前这一拍，不知道前面唱了什么，完全不会作词；
        \item RNN：有一本随手记录的小本子，每写一个字都翻翻本子，但本子页数有限，写得太长了前面的字迹会模糊；
        \item 古诗生成模型：拿着一本装满唐诗的课本练习"根据上文续写下一个字"，练多了，隐约学会了格律和意境——虽然还达不到李白的水平，但偶尔也能写出"月落乌啼霜满天"这样的意境。
    \end{itemize}
\end{metaphorbox}

\end{document}
